\subsection{Invertible submodules}

*(We preserve the notation of no. 5.)*

DEFINITION 4. A sub-A-module M of B is called invertible if there exists a sub-A-module N of B such that M.N = A.

Example. If \(b\) is invertible element of \(\mathbf{B}\) , the A-module \(Ab\) is invertible, as is seen by taking \(\mathbf{N} = \mathbf{Ab}^{-1}\) .

PROPOSITION 10. Let M be an invertible sub-A-module of B. Then:

\begin{enumerate}
\def\labelenumi{(\roman{enumi})}
\item
  There exists \(s \in S\) such that \(As \subset M \subset As^{-1}\) (and in particular \(M\) is nondegenerate).
\item
  A: M is the only sub-A-module N of B such that M.N = A.
\end{enumerate}

If M.N = A, then B.M = B.(B.M) ⊃ B.(M.N) = B.A = B, hence M is non-degenerate. Similarly N is non-degenerate. If \(t \in S \cap M\) and \(u \in S \cap N\) (no. 5, Proposition 8), the element s = tu belongs to \(S \cap M \cap N\) , whence \(Ms \subset M.N = A\) and therefore \(As \subset M \subset As^{-1}\) .

On the other hand, obviously \(N \subset A: M\) , whence

\[
\mathbf {A} = \mathbf {M}. \mathbf {N} \subset \mathbf {M}. (\mathbf {A}: \mathbf {M}) \subset \mathbf {A}
\]

and M. (A: M) = A; multiplying the two sides by N, we deduce A: M = N, which completes the proof.

THEOREM 4. Let M be a non-degenerate sub-A-module of B. The following properties are equivalent:

\begin{enumerate}
\def\labelenumi{(\alph{enumi})}
\item
  M is invertible.
\item
  M is projective.
\item
  M is projective of rank 1.
\item
  M is a finitely generated A-module and, for every maximal ideal m of A, the \(A_{m}\) -module \(M_{m}\) is monogenous.
\end{enumerate}

Let us show first the equivalence of properties (a), (b) and (c). If (a) holds and N is sub-A-module of B such that M.N = A, then there is a relation

\[
\sum_ {i = 1} ^ {p} m _ {i} n _ {i} = 1 \quad (m _ {i} \in \mathbf {M}, n _ {i} \in \mathbf {N} \text {   for   all   } i).\tag{9}
\]

\[
\boldsymbol {x} \in \mathbf {M},
\]

\[
v _ {i} (x) = n _ {i} x;
\]

\[
\nu_ {i}
\]

\(x = \sum_{i=1}^{n} m_i v_i(x)\) for all \(x \in \mathbf{M}\) ; this proves (Algebra, Chapter II, § 2, no. 6, Proposition 12) that \(\mathbf{M}\) is projective and generated by the \(m_i\) ; hence \(\mathbf{M}\) is a finitely generated projective module.

Let \(\mathfrak{m}\) be a maximal ideal of \(\mathbf{A}\) ; we show that the integer \(r = \operatorname{rg}_{\mathfrak{m}}(\mathbf{M})\) is equal to 1. Let \(S'\) be the image of \(S\) in \(A_{\mathfrak{m}}\) ; as the elements of \(S\) are not divisors of 0 in \(A\) , those of \(S'\) are not divisors of 0 in \(A_{\mathfrak{m}}\) , since \(A_{\mathfrak{m}}\) is a flat \(A\) -module (§2, no. 4, Theorem 1 and Chapter I, §2, no. 4, Proposition 3); then \(S'^{-1}A_{\mathfrak{m}} \neq 0\) and, as \(M_{\mathfrak{m}}\) is a free \(A_{\mathfrak{m}}\) -module of rank \(r\) , \(S'^{-1}M_{\mathfrak{m}}\) is a free \(S'^{-1}A_{\mathfrak{m}}\) -module of rank \(r\) . But if \(T'\) is the image of \(A - \mathfrak{m}\) in \(S^{-1}A\) , \(S'^{-1}A_{\mathfrak{m}}\) (resp. \(S'^{-1}M_{\mathfrak{m}}\) ) is canonically identified with \(T'^{-1}(S^{-1}A)\) (resp. \(T'^{-1}(S^{-1}M)\) ) (§2, no. 3, Proposition 7). Now \(S^{-1}M = B\) (Proposition 8 (c)) and hence \(T'^{-1}(S^{-1}M)\) is a free \(A\) -module of rank 1 over \(T'^{-1}(S^{-1}A)\) , which proves that \(r = 1\) and shows the implication (a) ⇒ (c).

The implication (c) \(\Rightarrow\) (b) is trivial. Let us show that (b) \(\Rightarrow\) (a). There exists by hypothesis a family (not necessarily finite) \((f_{\lambda})_{\lambda \in L}\) of linear forms on M and a family \((m_{\lambda})_{\lambda \in L}\) of elements of M such that, for all \(x \in M\) , the family \((f_{\lambda}(x))\) has finite support and \(x = \sum_{\lambda \in L} m_{\lambda} f_{\lambda}(x)\) (Algebra, Chapter II, §2, no. 6, Proposition 12). Since M is non-degenerate, \(f_{\lambda}(x) = n_{\lambda} x\) for some \(n_{\lambda} \in A: M\) by virtue of Proposition 9 of no. 5. Taking \(x\) as an element of M \(\cap S\) (no. 5, Proposition 8), it is seen that of necessity \(n_{\lambda} = 0\) except for a finite number of indices and \(\sum_{\lambda \in L} m_{\lambda} n_{\lambda} = 1\) . This obviously implies M. (A: M) = A, whence (a).

By virtue of Definition 2 of no. 3, (c) implies (d). Let us show the converse. As M is non-degenerate, its annihilator is zero (Proposition 8 (b)), then so is the annihilator of \(M_{m}\) (§ 2, no. 4, formula (9)). As \(M_{m}\) is assumed to be a monogenous \(A_{m}\) -module, it is therefore free of rank 1 and it then follows from no. 3, Theorem 2 that M is projective of rank 1.

COROLLARY. Every invertible sub-A-module of B is flat and finitely presented. This follows from Theorem 4 (c).

PROPOSITION 11. Let M, N be two sub-A-modules of B. Suppose that M is invertible. Then:

\begin{enumerate}
\def\labelenumi{(\roman{enumi})}
\tightlist
\item
  The canonical homomorphism \(\mathbf{M} \otimes_{\mathbf{A}} \mathbf{N} \to \mathbf{M}. \mathbf{N}\) is bijective.
\end{enumerate}

\[
\text {(ii)} \mathbf {N}: \mathbf {M} = \mathbf {N}. (\mathbf {A}: \mathbf {M}) a n d \mathbf {N} = (\mathbf {N}: \mathbf {M}). \mathbf {M}.
\]

Let \(j\) be the canonical injection \(N \to B\) . Since \(M\) is a flat A-module (Corollary to Theorem 4), \(1 \otimes j \colon M \otimes_{A} N \to M \otimes_{A} B\) is injective. But, as \(B = S^{-1}A\) , the B-module \(M \otimes_{A} B\) is equal to \(S^{-1}M\) and hence is identified with \(B\) since \(M\) is non-degenerate (no. 5, Proposition 8). If this identification is made, the image of \(1 \otimes j\) is \(M \cdot N\) , whence (i).

Let us set \(\mathbf{M}' = \mathbf{A} : \mathbf{M}\) . Then obviously \(\mathbf{M}' . \mathbf{N} \subset \mathbf{N} : \mathbf{M}\) and \(\mathbf{M} . (\mathbf{N} : \mathbf{M}) \subset \mathbf{N}\) . On the other hand, since \(\mathbf{M} . \mathbf{M}' = \mathbf{A}\) (Proposition 10),

\[
\mathbf {N}: \mathbf {M} = \mathbf {M} ^ {\prime}. \mathbf {M}. (\mathbf {N}: \mathbf {M}) \subset \mathbf {M} ^ {\prime}. \mathbf {N}
\]

and \(\mathbf{N} = \mathbf{M}.\mathbf{M}'.\mathbf{N}\subset \mathbf{M}.(\mathbf{N}:\mathbf{M})\) , whence (ii).

Remark. The proof of (i) in Proposition 11 uses only the fact that \(\mathbf{M}\) is flat and non-degenerate.

